% 与 example.tex 的版心兼容，给中文长行保留少量断行弹性。
\emergencystretch=3em
\section{问题二的建模与求解}
\subsection{问题二分析}
问题一把各服务区分开组批，因而未涉及实体无人机和电池在时间上的竞争。问题二允许同一架次访问多个服务区；一条路线的箱数、访问顺序和机型一经改变，沿途载荷、送达时刻及返航电量都会随之变化。即使若干架次分别满足航程约束，同时起飞也可能超过同型运输机或共享电池的库存。因此，本问须同时决定“哪些架次执行”和“何时由哪套设备执行”。

本文把单架次的空间与物理计算放在候选路线预处理中，把货箱覆盖、时限和资源周转放在整数排程中处理。前者保留多点访问的可能性，后者使运输机与电池的占用时段能够统一比较。图~\ref{fig:q2-resources}展示最终方案的两类资源时间轴；其中浅色区间是电池返航后充电及保守预留时间，不能作为下一架次的可用时间。

两问沿用问题一的航段时间和载荷相关能耗函数。为使本节可独立插入正文，约定$E_g^{\rm use}$为机型$g$的可用满电能量，$\rho_g$为返航最低电量比例，$T_g^{\rm full}$为题给从空电至满电的充电时长；$D_b,H_b,w_b$依次表示箱$b$的期望交付时刻、适用的硬截止时刻与优先系数。其中$E_g^{\rm use}$和$\rho_g$分别对应问题一记号$E_g$和$p_g$。后文所有任务开始时刻都指\emph{准备开始}，实际起飞还需等待准备和装载结束。这一区别影响交付偏移和资源时间轴。

返航电量比例记为$u=1-E/E_g^{\rm use}$。将题给前90\%与后10\%电量区间的充电时间分别线性分配，得到从比例$u$充至满电的时间
\begin{equation}
T_g^{\rm chg}(u)=
\begin{cases}
T_g^{\rm full}[0.65(0.90-u)/0.90+0.35],&u<0.90,\\
0.35T_g^{\rm full}(1-u)/0.10,&u\geq0.90.
\end{cases}
\label{eq:q2-charge}
\end{equation}
架次返航即释放运输机，所用电池在充满前仍占用库存；因此排程中必须分别追踪这两条时间线。

\subsection{基本约定与候选架次构造}
记$\mathcal B$为80个不可拆货箱集合，$\mathcal P$为预处理得到的候选运输架次。一个候选$p$包含货箱集合$B_p$、机型$g_p$和有序服务区序列$(a_1,\ldots,a_k)$，其路径为$O_{01}\to a_1\to\cdots\to a_k\to O_{01}$。同一服务区的货箱在一次停靠中交接；每个架次从满电池出发，到中心返航后才可更换或复用资源。货箱交付时刻取对应服务区卸载和交接结束时刻。

\subsubsection{逐航段载荷、时间和能源}

设$p$在航段$j$开始前尚未交付的货箱质量为$q_{pj}$，则$q_{p1}=\sum_{b\in B_p}m_b$。若第$j$站为$a_j$，离开该站后的载荷按$q_{p,j+1}=q_{pj}-\sum_{b\in B_p\cap B_{a_j}}m_b$更新；返航航段的载荷为零。因而跨区访问顺序不仅改变各箱送达时刻，也改变后续航段的能耗。

设架次$p$中第$j$个航段的起点和终点分别为$u_j$和$v_j$。沿用问题一建立的航段飞行时间模型$t_{ij}(g)$和航段能耗模型$E_{ij}(g\mid q)$，则架次$p$的总作业时间和总能耗分别为

\begin{equation}
\begin{aligned}
d_p={}&t^{\rm prep}_{g_p}+|B_p|t^{\rm load}_{g_p}
+\sum_{j\in p}t_{u_jv_j}(g_p)\\
&+\sum_{a\in p}
\left(
t^{\rm service}_{g_p}
+|B_p\cap B_a|t^{\rm unload}_{g_p}
\right),
\\
E_p={}&\sum_{j\in p}E_{u_jv_j}(g_p\mid q_{pj}).
\end{aligned}
\label{eq:q2-route-physics}
\end{equation}

对停靠序列中的服务区$a_k$，记到达该区的航段为$u_k\to a_k$，其中$u_1=O_{01}$、$u_j=a_{j-1}$（$j>1$）；令$B_{a_j}$为目的地是$a_j$的货箱集合。交接完成才算送达，因此属于$a_k$的箱$b$有
\begin{equation}
\delta_{pb}=t^{\rm prep}_{g_p}+|B_p|t^{\rm load}_{g_p}
+\sum_{j=1}^{k}t_{u_j a_j}(g_p)
+\sum_{j=1}^{k}\left(t^{\rm service}_{g_p}
+|B_p\cap B_{a_j}|t^{\rm unload}_{g_p}\right).
\label{eq:q2-delivery-offset}
\end{equation}
式~\eqref{eq:q2-delivery-offset}不含从最后一个服务区返航的时间，而式~\eqref{eq:q2-route-physics}中的$d_p$包含返航。因此单架次可以在返回中心之前完成多个不同时间的交付。对首批保障箱，用实际交接结束时刻判断时限，不能用起飞或到站时刻替代。

候选架次只有在起飞质量、体积、逐段航迹以及
\[
E_p\le(1-\rho_{g_p})E_{g_p}^{\rm use}
\]
均满足要求时才予以保留；即使立即开始准备也无法满足硬时限的候选同样剔除。这里的能量约束考察的是整个架次的返航余量，不能把每段分别可飞误当作整条路线可飞。

\subsubsection{有限候选池的构造}
候选池先纳入问题一的单区组批、每箱单独运输以及同区双箱组合，确保每个货箱至少存在基本投送选择。随后以固定随机种子进行55轮单区与55轮跨区受限组批：逐步增加货箱，对新服务区尝试不同插入位置，每条路线至多访问三个服务区；最后对已构造的货箱组合与访问顺序尝试其他可行机型。以“机型—货箱集合—访问顺序”去重，得到3988条候选，其中单区1783条、跨区2205条。完整货箱覆盖由后续式~\eqref{eq:q2-cover}保证；候选池覆盖所有箱并不意味着穷举了全部可能组批和访问顺序，因此不能从池内搜索结果推出原问题的全局最优。

\subsection{货箱覆盖、时限与共享资源模型}
\subsubsection{一次交付和时限}
令$x_p\in\{0,1\}$表示是否执行候选$p$，$s_p\in\mathbb Z_{\geq0}$为其准备开始时刻。任何货箱恰由一条选中路线交付：
\begin{equation}
\sum_{p:\,b\in B_p}x_p=1,\qquad b\in\mathcal B.
\label{eq:q2-cover}
\end{equation}
箱$b$在被选中的架次$p$中完成交接的实际时刻为$C_b=s_p+\delta_{pb}$。对具有硬截止时刻$H_b$的首批保障箱及医疗箱施加$C_b\le H_b$；题给期望时刻$D_b$用于衡量迟到，令$T_b=\max\{0,C_b-D_b\}$。当同一箱同时受到首批保障和医疗时限约束时，取较早的硬截止时刻。未设置硬截止的箱可以迟于$D_b$，但会付出加权迟到代价。
整数秒求解时，$s_p$为整数且各候选交付偏移向上取整，即$\widehat C_b=s_p+\lceil\delta_{pb}\rceil$。硬时限在求解器内写为$\widehat C_b\le\lfloor H_b\rfloor$；这样选中的排程回到连续物理时间后仍满足$C_b\le H_b$，但边界上可能排除本可行的非整数秒方案。该保守性是后文最优性结论只针对整数秒模型的原因之一。

\subsubsection{运输机与共享电池的双区间占用}
记机型$g$现有运输机数和电池数为$N_g^{\rm drone},N_g^{\rm bat}$，本题分别为A型$4/6$、B型$2/4$、C型$2/4$。运输机占用到实际返航，电池则还须经历式~\eqref{eq:q2-charge}给出的分段充电。为在整数秒排程中保守处理非整数时长，定义
\begin{equation}
\widehat d_p=\lceil d_p\rceil,\qquad
\widehat b_p=\left\lceil d_p+T_{g_p}^{\rm chg}
\left(1-\frac{E_p}{E_{g_p}^{\rm use}}\right)\right\rceil.
\label{eq:q2-reserved-length}
\end{equation}
选择架次$p$时，运输机占用$[s_p,s_p+\widehat d_p)$，电池占用$[s_p,s_p+\widehat b_p)$。对每种机型和任意时刻$t$，两类并发数分别不得超过其库存：
\begin{align}
\sum_{p:g_p=g}x_p\mathbf 1_{[s_p,s_p+\widehat d_p)}(t)
&\le N_g^{\rm drone},\\
\sum_{p:g_p=g}x_p\mathbf 1_{[s_p,s_p+\widehat b_p)}(t)
&\le N_g^{\rm bat}.
\label{eq:q2-cumulative}
\end{align}
整数模型中两组区间分别用可选区间与累积容量约束表达。求解后再按区间不重叠原则把任务匹配至具体机号、电池号，并以未取整的物理时刻复核。此时表~\ref{tab:q2-result}中的完工时间按实际返航计算，与整数目标中的向上取整完工值略有差别。

这里把$H_b$与$D_b$分开是必要的：前者一旦违反，方案即不可行；后者只改变方案的优劣。对没有硬截止的箱不人为补造截止时刻；对两项硬要求同时适用的箱取较早时刻，也避免重复计算。半开资源区间允许上一任务的释放时刻恰好等于下一任务的开始时刻，与充电完成后立即复用电池一致。未选路线的可选区间不占用库存，因此货箱覆盖和资源容量可在同一模型内联立求解。

表~\ref{tab:q2-model-map}将候选预处理与排程求解的边界列出。航程、载荷和返航电量在生成候选时已检查，因而整数模型无须重复展开每个飞行航段；相反，同型机和电池的同时占用取决于所有已选架次的起始时刻，只能在排程层联立判断。这种划分减少了组合模型的规模，也使回代时可以分别核查“路线自身是否可飞”和“多条路线能否同时执行”。
\begin{table}[htbp]
\centering\small
\caption{问题二两层计算中的主要变量、约束与核查对象}
\label{tab:q2-model-map}
\begin{tabular}{@{}p{0.17\textwidth}p{0.36\textwidth}p{0.39\textwidth}@{}}
\toprule
层次 & 决策或计算量 & 可行性条件与作用\\
\midrule
路线预处理 & 货箱集合、机型、停靠顺序；$d_p,E_p,\delta_{pb}$ & 质量、体积、逐段载荷及返航电量；剔除即时起飞亦不满足硬时限的路线\\
整数排程 & 路线选择$x_p$、准备开始$s_p$、整数交付$\widehat C_b$ & 80箱恰好覆盖一次；硬截止时刻；运输机与电池分别满足累积容量\\
物理回代 & 实际交付$C_b$、返航时刻与组件就绪时刻 & 恢复未取整航段值，逐箱、逐资源检查最终排程\\
\bottomrule
\end{tabular}
\end{table}

\subsection{优化目标与求解方法}
\subsubsection{整数化加权目标}
本问同时关心送达及时性、最后返航时间、能耗和架次数。为使求解器能够直接比较不同候选路线，采用代码中\texttt{balanced}配置的单一整数目标。设$\widehat C_b=s_p+\lceil\delta_{pb}\rceil$、$\widehat T_b=\max(0,\widehat C_b-D_b)$，并令$\widehat M=\max_{p:x_p=1}(s_p+\widehat d_p)$，则
\begin{equation}
\min\;100\widehat M
+\sum_{b\in\mathcal B}w_b\bigl(\widehat C_b+100\widehat T_b\bigr)
+\sum_{p\in\mathcal P}x_p\left[10\,\operatorname{round}(1000E_p)+10000\right].
\label{eq:q2-objective}
\end{equation}
这里$w_b$是题给优先系数。式~\eqref{eq:q2-objective}是加权和，并非严格字典序；报告“迟到为零”是求得排程的事实，不应理解为目标函数强制零迟到。能耗项以Wh取整，所有系数均为本次计算采用的权重，不把它们解释成真实货币成本。

式中的$\widehat M$由已选架次的整数返航时刻最大值等式实现，$\widehat T_b$由非负变量及$\widehat T_b\geq\widehat C_b-D_b$实现。迟到项系数为正，故最优时$\widehat T_b$取可行最小值。求解器中的秒级向上取整用于避免低估资源占用，正式结果仍以未取整的航段物理量回代。

\subsubsection{候选池求解与固定路线重排}
CP-SAT先在有限候选池内联合选择路线和任务开始时刻，原始排程目标值为4290638。随后固定24条选中路线和货箱分配，只重排整数秒开始时刻及两类容量资源。新目标值为4281382，较原排程减少9256（约0.216\%）。

固定路线子问题的求解状态为\texttt{OPTIMAL}，目标值与下界同为4281382。

表~\ref{tab:q2-objective-parts}按式~\eqref{eq:q2-objective}拆开重排前后的目标值。路线和货箱分配未变，故架次数与能耗项不变；联合返航完工项也相同。9256的改善全部来自货箱加权交付时刻提前，而不是架次数或能耗减少。这个拆解同时检验了目标函数计算与正式逐箱交付表的一致性。
\begin{table}[htbp]
\centering\small
\caption{固定路线重排前后目标值分项；数值均为式~\eqref{eq:q2-objective}的整数分值}
\label{tab:q2-objective-parts}
\begin{tabular}{@{}lrrr@{}}
\toprule
分项 & 原排程 & 重排后 & 改变量\\
\midrule
联合完工项 & 699800 & 699800 & 0\\
加权交付项 & 2655038 & 2645782 & $-9256$\\
加权迟到项 & 0 & 0 & 0\\
能耗项 & 695800 & 695800 & 0\\
架次数项 & 240000 & 240000 & 0\\
\midrule
总目标值 & 4290638 & 4281382 & $-9256$\\
\bottomrule
\end{tabular}
\end{table}

为考察更换组批和访问顺序的可能改进，还进行了20轮大邻域重建：每轮移出5--9条当前路线，补入相关替代路线和部分成对合并路线，给CP-SAT 8\,s搜索时间，随机种子为20260923。20轮均未接受更优解；因此，本次已实现的目标值改善来自前述\emph{固定路线重排}，不能归因于大邻域搜索。“固定路线最优”也仅针对这24条路线、整数秒开始时刻和给定资源模型，不证明3988条候选池或全部可行路线上的全局最优。

\subsection{计算结果与分析}
\subsubsection{运输任务与机型分工}
表~\ref{tab:q2-result}给出正式方案。24个架次交付80箱，A、B、C型分别执行10、8、6架次；其中C型6架次承运39箱，体现大载重机型的组批能力。仅两条架次访问两个服务区：\texttt{Q2-009}沿$S009\to S012$运行，\texttt{Q2-020}沿$S011\to S014$运行。跨区比例较小，说明在给定地形、时限和能源条件下，多点访问并非对每批货箱都合算。

作为候选策略对照，把问题一的18条既定组批路线原样施加问题二的时限和双资源约束，模型返回不可行；仅从1783条单区候选中求得的可行排程为26架次、71.599\,kWh。允许2205条跨区候选参加选择后，找到24架次、69.580\,kWh的方案，随后固定这24条路线优化时序。表~\ref{tab:q2-strategy}把这些\emph{已求得}的方案放在同一物理指标下比较。单区方案和跨区方案的求解均受候选池与时间限制，因此不能据此证明跨区路线在所有可行方案中必然节省2架次。
\begin{table}[htbp]
\centering\small
\caption{问题二不同候选策略下求得的运输方案}
\label{tab:q2-strategy}
\begin{tabular}{@{}lrrrr@{}}
\toprule
方案 & 架次 & 跨区架次 & 电耗/kWh & 加权平均交付/s\\
\midrule
仅单区候选 & 26 & 0 & 71.599 & 2781.010\\
跨区候选原排程 & 24 & 2 & 69.580 & 2776.663\\
固定路线重排后 & 24 & 2 & 69.580 & 2766.981\\
\bottomrule
\end{tabular}
\end{table}

\begin{table}[htbp]
\centering\small
\caption{问题二正式运输方案的汇总结果}
\label{tab:q2-result}
\begin{tabular}{@{}lrrr@{}}
\toprule
指标 & A型 & B型 & C型\\
\midrule
运输架次 & 10 & 8 & 6\\
交付货箱数 & 20 & 21 & 39\\
运输能耗/kWh & 27.125 & 17.716 & 24.740\\
\midrule
\multicolumn{3}{@{}l}{全部货箱/全部架次} & 80 / 24\\
\multicolumn{3}{@{}l}{加权迟到/迟到箱数} & 0 / 0\\
\multicolumn{3}{@{}l}{运输返航完工时间/s} & 6997.523\\
\multicolumn{3}{@{}l}{总运输能耗/kWh} & 69.580\\
\multicolumn{3}{@{}l}{最小返航电量比例} & 23.088\%\\
\multicolumn{3}{@{}l}{最小硬时限剩余时间/s} & 220.525\\
\bottomrule
\end{tabular}
\end{table}

图~\ref{fig:q2-resources}按具体机号和电池号绘制排程。前8个架次在任务起点同时开始准备，之后的任务由已返航运输机及可用电池接续执行。由正式区间直接求得A、B、C型运输机峰值分别为$4/2/2$架，共享电池峰值分别为$6/4/4$块，恰与各型现有库存相等。若只看运输机的空闲而忽略下方浅色充电区间，就会高估可同时出动的能力；本图也给出了同型电池在不同实体机之间周转的实际安排。
\begin{figure}[htbp]
\centering
\includegraphics[width=0.96\textwidth]{figures/result_q2_resources.pdf}
\caption{问题二正式方案的运输机与共享电池占用时间轴；浅色部分为返航后的电池充电和保守预留区间}
\label{fig:q2-resources}
\end{figure}

\subsubsection{送达时效与安全余量}
80箱均在期望时刻前交付，加权迟到为零；按优先系数加权的平均交付时刻为2766.981\,s。具有硬截止时刻的货箱在图~\ref{fig:q2-slack}中均位于零线右侧。最紧的是饮用水箱\texttt{S012-\allowbreak{}WAT-01}，仍提前220.525\,s。若临时延迟相关架次，应首先复查此箱的硬时限。所有架次最小返航电量比例为23.088\%。
\begin{figure}[htbp]
\centering
\includegraphics[width=0.86\textwidth]{figures/result_q2_hard_deadline_margin.pdf}
\caption{问题二各硬时限货箱的交付余量}
\label{fig:q2-slack}
\end{figure}

\subsection{正确性与最优性范围检验}
正式结果保存了逐架次、逐箱、逐航段和电池周转记录。逐箱核对80个箱的唯一交付、交付时刻与对应服务区；逐航段重新累加飞行时间和载荷相关能耗，检查起飞质量、体积及返航电量；再将每条路线配给具体机号、电池号，检查同一资源的占用区间是否相交。表~\ref{tab:q2-audit}列出相互独立的核查对象与结果。硬时限和期望时刻仍按各自定义分别计算，不能用“迟到箱数为零”代替硬约束验证。
\begin{table}[htbp]
\centering\small
\caption{问题二正式排程的逐项回代核查}
\label{tab:q2-audit}
\begin{tabular}{@{}p{0.26\textwidth}p{0.47\textwidth}r@{}}
\toprule
核查对象 & 回代依据 & 结果\\
\midrule
不可拆货箱 & 逐箱记录与80箱清单逐一匹配 & 80箱唯一交付\\
硬截止时刻 & 对适用的箱计算$H_b-C_b$ & 最小余量220.525\,s\\
返航电量 & 逐航段电耗除以机型可用电能 & 最小23.088\%\\
运输机与电池 & 具体设备占用、返航与充满区间 & 冲突0次\\
固定路线整数排程 & 目标值与该子问题下界 & 均为4281382\\
\bottomrule
\end{tabular}
\end{table}

最后一行的最优证书只证明已选24条路线在给定整数秒模型下的排程最优。候选生成没有穷举所有货箱组合，原候选池求解也没有匹配全局上下界；20轮大邻域搜索未接受新解只能说明这些受限邻域没有找到改进。故问题二的完整结果应称为\emph{物理与资源约束均已验证的可行调度}。第三问将加入连续通信约束，并允许调整部分运输时刻。

\section{问题三的建模与求解}
\subsection{问题三分析}
问题二的24架次排程满足运输与电池周转约束，却没有检查飞行过程中能否持续与固定网关$G_{01}$通信。山区遮挡使“起飞点、服务区两端可连通”不能推出整条航迹可连通；中继又占用两架中继机和六个能源组件，其建链、悬停及充电会反过来限制运输任务开始时刻。本问因而在已有货箱分配的基础上联动调整运输时序、中继位置和服务窗口，要求每一条运输轨迹的全部需通信时段都存在有效双向链路。

初始第二问排程的直连失效比例在不同架次间差异明显，见图~\ref{fig:q3-direct-gap}。图中比例是预计算网格上的诊断量，并非最终连续通信证明。中继位置若只为一个架次选择，可能产生大量短窗口；若能覆盖多条轨迹，就有机会减少中继起降和组件占用。因此，求解中把覆盖能力与任务时序同时考虑。
\begin{figure}[htbp]
\centering
\includegraphics[width=0.82\textwidth]{figures/raw_q3_q2_direct_outages.pdf}
\caption{问题二初始运输架次的网格化直连失效比例；通信约束尚未施加}
\label{fig:q3-direct-gap}
\end{figure}

\subsection{通信链路与连续覆盖模型}
通信要求作用于运输机需要保持联络的实际轨迹时间，而非仅作用于服务区或航段端点。将爬升、巡航、下降、交接及返航拆成分段轨迹；中继机到达悬停点并完成建链后才成为可用节点，返航途中不再提供覆盖。这样，通信可行性成为时间与空间同时成立的约束。
\subsubsection{双向链路预算}
对端点$a,b$，由发射功率、收发增益、系统损耗和接收灵敏度计算两个方向的允许路径损耗，并取较小值作为双向门槛。将题给8\,dB通信裕量加入各接收灵敏度，定义有效接收门槛$P_{{\rm th},i}=P_{{\rm sens},i}+8\,\mathrm{dB}$，则
\begin{equation}
L^{\max}_{a\leftrightarrow b}
=\min\!\left\{
P_a+G_a^t+G_b^r-L_{\rm sys}-P_{{\rm th},b},\;
P_b+G_b^t+G_a^r-L_{\rm sys}-P_{{\rm th},a}
\right\}.
\label{eq:q3-budget}
\end{equation}
式中发射功率与接收门槛以dBm计，天线增益以dBi计，系统损耗与路径损耗以dB计。给定时刻$t$两端点三维距离$D_{ab}(t)$（km）、工作频率$f$（MHz），用DEM判定视线是否被地形遮挡，记遮挡指示量为$z_{ab}(t)$，则
\begin{equation}
L_{ab}(t)=32.45+20\log_{10}f+20\log_{10}D_{ab}(t)
+L_{\rm obs}z_{ab}(t),\qquad
A_{ab}(t)=\mathbf1\!\left\{L_{ab}(t)\le L^{\max}_{a\leftrightarrow b}\right\}.
\label{eq:q3-link}
\end{equation}
这里$L_{\rm obs}$是题给遮挡附加损耗。运输机位置按准备后爬升、水平巡航、下降和服务区交接阶段分段计算，而非把整条路线视为一根静态直线。

\subsubsection{直连或单中继两跳保障}
记运输机$p$的通信需求时间集合为$\mathcal T_p$。若中继任务$r$已完成建链并处于悬停服务窗口，令$I_r(t)=1$；否则为0。对所有$p$及$t\in\mathcal T_p$，要求
\begin{equation}
A_{pG}(t)=1
\quad\text{或}\quad
\exists r:\ I_r(t)A_{pr}(t)A_{rG}(t)=1.
\label{eq:q3-cover}
\end{equation}
因此，两跳链路必须在\emph{同一时刻}同时成立，不能分别用不同时刻的可达性代替。模型只使用固定网关直连或“运输机—中继—网关”两跳，不引入多中继转发。题目没有给出单架中继的并发接入上限；在这一数据边界下，同一中继服务窗可同时保障多架运输机。

式~\eqref{eq:q3-cover}等价于对每个$t\in\mathcal T_p$要求至少有一个完整的通信提供者。提供者可以随时间切换，不必强迫一条运输路线始终依赖同一架中继，但切换处也必须有可用链路。只在固定步长的样点上检查，无法排除两样点之间的短暂遮挡；后续连续证书专门处理这一缺口。

\subsection{中继任务与能源资源模型}
中继并非可随意放置的静态覆盖点，而是一条从中心出发、建立回传链路、悬停服务、再返回中心的完整架次。候选点只是搜索空间；决定采用某点时，服务开始、结束、机号和能源组件仍需与运输任务共同排程。长服务窗有利于共享保障，却会增加悬停电耗并延后同一中继机或组件的下次可用时刻。
\subsubsection{候选悬停点与预筛}
候选点以服务区、中心至服务区航线上的四分点和中点及服务区附近点为基础，并在部分位置枚举不同离地高度。每个位置均由DEM取得地面高程，悬停高度不超过题给上限。候选点需同时满足中继至网关的双向回传链路、往返飞行能量和返航余量。对每条运输轨迹上的直连缺口，预计算可由哪些候选点接入；图~\ref{fig:q3-sites}列出覆盖初始失联采样较多的一组候选点，说明位置选择直接影响后续可合并的服务窗口数量。该图用于筛选和解释搜索，不作为最终连续可行性的证据。
\begin{figure}[htbp]
\centering
\includegraphics[width=0.82\textwidth]{figures/process_q3_candidate_coverage.pdf}
\caption{候选中继点对初始运输轨迹直连失效采样点的覆盖数量；每个候选点独立统计}
\label{fig:q3-sites}
\end{figure}

候选覆盖采用粗轨迹采样，目的是淘汰明显无效的位置，而非证明正式方案。束搜索采用候选位置的若干高度层，先检查回传和往返飞行余量，再按10\,s网格记录运输机与候选点的可达关系。若某点能覆盖多条路线的直连缺口，可通过延长悬停窗复用一次中继；若所需窗口过长或能耗超限，必须拆分架次或换点。

\subsubsection{一次中继架次的时间与电量}
一条中继架次从$O_{01}$起飞，经历准备、往程、建链、悬停服务、返程和周转。设候选点为$k_r$，建链后服务窗口为$[\sigma_r,\tau_r]$，则其电耗写为
\begin{equation}
E_r^{\rm relay}=E_r^{\rm flight}
+\frac{P_{\rm hover}+P_{\rm comm}}{3600}
\bigl(T_{\rm link}+\tau_r-\sigma_r\bigr).
\label{eq:q3-relay-energy}
\end{equation}
$E_r^{\rm flight}$由往返水平飞行功率和爬升势能共同计算；$P_{\rm hover},P_{\rm comm}$以kW计，式中时间以s计，因此除以3600后得到kWh。服务项含建链时间，避免漏计建链期间的悬停与通信耗电。中继电池返航后须保有规定比例电量。两架实体中继机的任务区间加上周转时间不得重叠；六个能源组件在其所执行任务返航并完成分段充电前不得复用。可行候选服务窗还须覆盖需要它的运输机失联时段，并在窗口两端留出预设时间缓冲。

服务窗$[\sigma_r,\tau_r]$只是实际任务的一部分。中继机占用从准备开始延续到返航及规定周转结束；组件占用则延续到返航充满，二者释放时刻通常不同。式~\eqref{eq:q3-relay-energy}中的服务时长增加，会同时推高电耗与组件充电时间，因此不能仅根据覆盖点的几何可见性决定中继数量。

具体地，令$\ell_r$为中继任务开始准备的时刻，$e_r$为返航时刻，$T^{\rm chg}_r$为根据返航电量计算的组件充电时间。中继机和能源组件分别占用半开区间
\begin{equation}
I_r^{\rm relay}=[\ell_r,e_r+300),\qquad
I_r^{\rm component}=[\ell_r,e_r+T^{\rm chg}_r).
\label{eq:q3-resource-interval}
\end{equation}
300\,s是题给的中继机返航周转时间。对同一机号和同一组件号，式~\eqref{eq:q3-resource-interval}中的相应区间不得相交；若不预先指定编号，则任意时刻的并发数分别不得超过2架中继机和6个能源组件。两种资源的释放顺序未必相同，例如中继机完成周转后可换用其他已就绪组件再次出动。实际排程同时检查这两类区间，以及中继往返飞行、建链和服务窗之间的先后关系。

\subsection{联合求解与时间细化}
若在连续时间内直接联立优化24条运输路线的时序、全部中继位置及服务窗，变量耦合很强。本文采用“粗网格生成候选—有限宽度联合排程—实例时间细化—精确回代与连续证明”的分层流程。粗网格只用于搜索，最终可行性由原始DEM、完整轨迹及时间区间证书决定。
表~\ref{tab:q3-fidelity}说明各阶段的时间分辨率和判定职责。候选点展示图~\ref{fig:q3-sites}仅使用全高度候选、20\,s初始运输轨迹样点；实际束搜索重新生成0.5、0.75和1倍高度的候选位置，并用10\,s轨迹网格建立覆盖关系。两者服务于不同步骤，不能把图中的覆盖数量当作最终服务时长或链路证明。
\begin{table}[htbp]
\centering\small
\caption{问题三搜索、回代和证明的分辨率与职责}
\label{tab:q3-fidelity}
\begin{tabular}{@{}p{0.20\textwidth}p{0.26\textwidth}p{0.46\textwidth}@{}}
\toprule
阶段 & 位置及时间处理 & 可支持的结论\\
\midrule
候选点展示 & 全高度位置，20\,s样点 & 比较初始失联样点的候选覆盖数，用于解释选点\\
束搜索 & 三层高度，10\,s轨迹网格 & 筛选服务窗与运输时序，形成待严格验证的候选\\
事件采样回代 & 完整轨迹、最大间隔1\,s，增补服务窗边界 & 检出样点处的失联及物理、时间、资源违反\\
连续区间证书 & 阶段端点起始、最长30\,s初分，未决时二分 & 在DEM与链路模型内证明所有需求时刻都有可用提供者\\
\bottomrule
\end{tabular}
\end{table}
\subsubsection{搜索状态与排序}
以问题二的24个架次和原有货箱分配为输入。为使具有3600\,s硬截止时刻的医疗箱\texttt{S012-MED-01}先交接，跨区架次\texttt{Q2-009}的访问顺序由$S009\to S012$调整为$S012\to S009$，并重新计算逐段载荷、航时、能耗、交付偏移和电池就绪时刻。表~\ref{tab:q3-order}比较的是\emph{从任务准备开始到交付的偏移}，剔除了两问任务开始时刻不同的影响。前往$S012$的交付偏移缩短295.234\,s，但$S009$的食品箱延后715.711\,s，路线能耗增加0.040388\,kWh；因此倒序不是无代价的改进。除此之外，搜索保留原有货箱分配与机型，主要调整任务时序和中继服务窗。
\begin{table}[htbp]
\centering\small
\caption{架次\texttt{Q2-009}改变访问顺序后的箱级交付偏移}
\label{tab:q3-order}
\begin{tabular}{@{}lrrr@{}}
\toprule
货箱 & 原顺序偏移/s & 倒序偏移/s & 变化/s\\
\midrule
\texttt{S012-MED-01} & 1599.709 & 1304.475 & $-295.234$\\
\texttt{S009-FOD-01} & 1095.897 & 1811.608 & $+715.711$\\
\bottomrule
\end{tabular}
\end{table}

先用10\,s网格标记各运输轨迹的网关直连缺口及候选中继点的两跳覆盖关系。问题二已选24个架次的承运货箱集合、机型及具体运输机、电池编号保持不变，只有上段所述\texttt{Q2-009}访问顺序改变；前8个架次固定在$t=0$同时开始准备，并据其直连缺口枚举初始中继点对。每个状态记录已排和未排架次、设备就绪时刻及中继任务的点位与服务窗。扩展待排架次时，依次尝试直连、复用或延长已有中继窗以及新建中继任务，并检查硬时限、能量和资源约束。排序同时考虑已发生迟到、未来紧急路线的最早可能迟到以及候选点稀缺程度；对同一初始点对及新中继点限制近似重复状态，每层最多保留160个。束宽和上述固定决策限制了搜索空间，算法不提供全局最优界。

本次实例枚举78组初始中继点对，形成79个初始搜索状态；随后对余下16条运输路线逐层扩展，累计生成37275个子状态，搜索耗时47.226\,s。表~\ref{tab:q3-search-flow}把每一层的判断顺序压缩为可复核的过程。排序分数中的“未来最早可能迟到”是启发式评估，并不是原问题目标函数的有效全局下界；同样，保留160个状态意味着被剪去的方案仍可能更优。搜索规模用于说明所实施的计算，而不构成穷举证明。
\begin{table}[htbp]
\centering\small
\caption{问题三有限宽度联合排程的状态扩展过程}
\label{tab:q3-search-flow}
\begin{tabular}{@{}p{0.15\textwidth}p{0.75\textwidth}@{}}
\toprule
步骤 & 处理与判定\\
\midrule
初始化 & 固定最先准备的8条运输路线；从其直连缺口构造中继点对和服务窗\\
扩展 & 对一条未排运输路线枚举可用开始时刻，并尝试网关直连、已有中继窗延长或新建中继任务\\
过滤 & 同时检查箱级硬时限、两类运输资源、中继能量以及式~\eqref{eq:q3-resource-interval}的双资源占用\\
截断 & 按迟到与未来紧急程度等排序，去掉近似重复状态，每层至多保留160个\\
验收 & 对完整状态作事件采样、连续区间证书及逐项物理回代，比较通过验证的候选\\
\bottomrule
\end{tabular}
\end{table}

完整候选以原始DEM和精确视线判定执行最大间隔1\,s的事件采样复核，再用连续时间证书检查样点之间的全部区间。经验证的候选按“加权迟到、联合完工时刻、总能耗、总架次数”字典序比较。束搜索中间方案使用9个中继架次，13箱迟到，加权迟到499378.079。随后针对本实例9个指定运输架次调整开始时刻，延长一条长时中继窗并重设最后一个中继任务，得到5个中继架次的正式方案；这属于实例时间细化，并非通用自动优化器。调整后的运输与中继任务重新接受物理、时限、资源和通信检验。

\subsubsection{连续时间通信证书}
逐秒检查只覆盖采样时刻，不能排除两采样点之间的短暂掉线。为消除这一缺口，按运输轨迹的运动阶段端点、中继服务窗端点及需要时的递归二分点划分时间。每个小区间内运输机位置是时间的仿射函数；对固定网关或悬停中继，整段视线的空间扫掠区域可用端点构成的三角形界定。将该区域与DEM像元逐一求交，保守确定整段是否有遮挡；同时对区间内最大三维距离取上界，代入式~\eqref{eq:q3-link}得到链路裕量下界。若直连不能在整段得到证明，再检查是否有同一中继两段链路均可整段证明；仍未决时继续二分。

对时间区间$I=[t_0,t_1]$，取运输机位置端点$u(t_0),u(t_1)$与固定提供者$v$，构造保守三维距离上界$\overline D(I)$；由式~\eqref{eq:q3-link}得到自由空间损耗上界$\overline L_0(I)$。若
\begin{equation}
L^{\max}_{u\leftrightarrow v}-\overline L_0(I)-L_{\rm obs}\geq\varepsilon_L,
\label{eq:q3-worst-block}
\end{equation}
则即使全区间按遮挡损耗计算仍可连通，是否存在视线遮挡无需再判定。否则，若无遮挡损耗也超出预算，就否定该提供者的整段证明；剩余情形才检查视线扫掠区域。设该运动阶段内$u(t)$随$t$仿射变化，连接运输机与固定提供者的射线可写为
\begin{equation}
q(t,\lambda)=(1-\lambda)u(t)+\lambda v,
\qquad t\in I,\quad 0\le\lambda\le1.
\label{eq:q3-swept-ray}
\end{equation}
式~\eqref{eq:q3-swept-ray}的所有点落在$u(t_0),u(t_1),v$张成的扫掠三角形内。对给定DEM像元，把此三角形裁切为多边形。题给DEM按像元常高处理，裁切多边形上射线高度的最小值必在某个顶点取得；若每个相交像元的最低顶点都高于该像元地面高程及高度保护量，便证明整个时间区间视线无阻。这一判定检查的是三角形内部所有射线，不是只检查区间两端。对中继提供者还须保证同一服务窗覆盖整个$I$，且悬停点到网关的双向回传链路成立。$\varepsilon_L$以及距离和高度保护量用于保守处理数值边界。

区间端点先包含运动阶段、中继服务窗边界，再按不超过30\,s的间距初分；某段找不到可证明的提供者时递归二分，直到证明成立或达到最小细分尺度。证明对象是区间内\emph{所有}时刻的链路预算和DEM视线，而非增加一次中点抽样。若存在未决区间，就拒绝该候选方案。

正式证书把全部36198.242\,s运输通信需求划为1317个已证区间，未决区间为零。这里的需求时间是各架次首段起飞至返航的时间之和，含途中交接，不含起飞前的准备和装载。其中966个区间使用式~\eqref{eq:q3-worst-block}的遮挡最坏情形界，351个区间使用扫掠DEM净空判定；后一类判定累计执行194870次DEM像元检查，同一像元在不同区间可重复出现。图~\ref{fig:q3-cert}展示两类区间的数量。最小保守链路裕量下界为0.008976\,dB，出现在\texttt{Q2-009}于3526.607--3533.489\,s经\texttt{R02}在\texttt{NS013}中继的区间；它仍由遮挡最坏情形界证明。证书证明的是题给DEM、分段仿射轨迹及既定链路预算下的连续可行性，不能将这一很小的模型裕量解释为足够的实地环境扰动余量。
\begin{figure}[htbp]
\centering
\includegraphics[width=0.73\textwidth]{figures/process_q3_certificate_methods.pdf}
\caption{连续通信证书中两种区间证明方法的使用次数}
\label{fig:q3-cert}
\end{figure}

\subsection{联合调度结果与分析}
\subsubsection{中间方案与正式方案}
表~\ref{tab:q3-result}把问题二运输基线、束搜索中间解和最终验证方案并列。第二问没有通信约束，因此其零迟到只能作为运输基线，不是第三问的可行目标值。正式方案仍用24个运输架次交付全部80箱，额外执行5个中继架次；联合完工时间8772.339\,s，以最后完成的运输或中继返航为准。总能耗73.835\,kWh，其中中继任务消耗4.215\,kWh。
\begin{table}[htbp]
\centering\small
\caption{问题二基线与问题三搜索结果比较}
\label{tab:q3-result}
\begin{tabular}{@{}lrrr@{}}
\toprule
指标 & 问题二基线 & 束搜索中间解 & 正式验证方案\\
\midrule
运输架次 & 24 & 24 & 24\\
中继架次 & 0 & 9 & 5\\
迟到箱数 & 0 & 13 & 1\\
加权迟到/优先系数$\cdot$s & 0 & 499378.079 & 8546.904\\
联合完工时间/s & 6997.523 & 13262.942 & 8772.339\\
运输能耗/kWh & 69.580 & --- & 69.621\\
中继能耗/kWh & --- & --- & 4.215\\
总能耗/kWh & 69.580 & 75.309 & 73.835\\
硬时限违反/通信中断 & 未检通信 & 0 / 0 & 0 / 0\\
\bottomrule
\end{tabular}
\end{table}

图~\ref{fig:q3-timeline}显示两类任务相互穿插：长时中继窗在较长时间内保持服务，其他四次中继集中补齐特定时段的直连缺口。表~\ref{tab:q3-relay}给出具体中继机、组件、位置和建链后的服务窗口；例如\texttt{R02}先后执行4次任务，但相邻任务的返航、周转和组件就绪要求仍逐项满足。
\begin{figure}[htbp]
\centering
\includegraphics[width=0.94\textwidth]{figures/result_q3_joint_timeline.pdf}
\caption{第三问运输与中继任务的联合时间安排；$T1$--$T24$为运输架次，$R1$--$R5$为中继架次}
\label{fig:q3-timeline}
\end{figure}
\begin{table}[htbp]
\centering\small
\caption{第三问五个中继架次的建链后服务窗口}
\label{tab:q3-relay}
\begin{tabular}{@{}llllrr@{}}
\toprule
架次 & 中继机 & 组件 & 悬停点 & 开始/s & 结束/s\\
\midrule
1 & R01 & RE01 & LS015-0.75 & 827.593 & 7169.953\\
2 & R02 & RE02 & S010-H50 & 732.084 & 1484.108\\
3 & R02 & RE03 & NS013 & 2887.432 & 3786.714\\
4 & R02 & RE02 & LS004-0.75 & 5305.070 & 6143.617\\
5 & R02 & RE04 & S010-H50 & 7623.693 & 8266.037\\
\bottomrule
\end{tabular}
\end{table}

\subsubsection{交付代价与通信保障}
表~\ref{tab:q3-relay-support}按连续通信证书统计各次中继实际承担的运输通信时间。首个中继任务的服务窗长6342.360\,s，却累计保障16个运输架次的9784.861\,s通信时间；由于同一时刻允许服务多架运输机，两者并不矛盾。其他四次中继主要填补长窗外或其覆盖不到的局部缺口。支持架次数按不同运输架次计，支持时长按各运输架次的已证区间求和，不能把两个数直接相除解释为平均并发量。
\begin{table}[htbp]
\centering\small
\caption{问题三中继任务的经证书确认的保障量与能耗}
\label{tab:q3-relay-support}
\begin{tabular}{@{}lrrrr@{}}
\toprule
任务 & 保障运输架次 & 累计保障/s & 电耗/kWh & 返航电量/\%\\
\midrule
1：\texttt{LS015-0.75} & 16 & 9784.861 & 2.199 & 31.289\\
2：\texttt{S010-H50} & 4 & 1625.853 & 0.473 & 85.209\\
3：\texttt{NS013} & 4 & 2557.850 & 0.546 & 82.935\\
4：\texttt{LS004-0.75} & 1 & 446.982 & 0.557 & 82.608\\
5：\texttt{S010-H50} & 1 & 642.345 & 0.440 & 86.256\\
\bottomrule
\end{tabular}
\end{table}

逐箱结果如图~\ref{fig:q3-box}。80箱均唯一交付，仅食品箱\texttt{S014-\allowbreak{}FOD-01}晚于期望时刻854.690\,s。其优先系数为10，形成8546.904优先系数秒的加权迟到，但不违反硬时限；其余箱均按期望时刻送达。

表~\ref{tab:q3-box-critical}把最紧的硬时限箱与唯一迟到箱放到问题二基线上比较。\texttt{Q2-012}的准备开始时刻推迟550.684\,s，使\texttt{S002}的医疗和饮用水箱恰在3600\,s硬截止时刻完成交接；其问题二原有的550.684\,s余量全部耗尽。\texttt{S014-FOD-01}所属\texttt{Q2-020}的准备开始由3598\,s推迟到6175\,s，交付时刻相应推迟2577\,s，从原先提前1722.310\,s变为迟到854.690\,s。后者是本次联合排程观察到的代价，不能单凭两份解的差值断言它是通信约束不可避免的最小代价，也不能把延迟归因于某一次中继任务。
\begin{table}[htbp]
\centering\small
\caption{问题二与问题三关键货箱的时刻变化}
\label{tab:q3-box-critical}
\begin{tabular}{@{}lrrrr@{}}
\toprule
货箱 & 时限或期望/s & 问题二送达/s & 问题三送达/s & 问题三余量或迟到/s\\
\midrule
\texttt{S002-MED-01} & 3600 & 3049.316 & 3600.000 & 硬余量0\\
\texttt{S002-WAT-01} & 3600 & 3049.316 & 3600.000 & 硬余量0\\
\texttt{S014-FOD-01} & 7200 & 5477.690 & 8054.690 & 迟到854.690\\
\bottomrule
\end{tabular}
\end{table}
因此，问题三虽然保持了所有硬约束可行，但最小硬时限余量由问题二的220.525\,s降为0。若需应对出发延迟或参数误差，应先为\texttt{S002}两箱重新留出时间余量；当前证书只验证了给定参数下的正式时刻表。
\begin{figure}[htbp]
\centering
\includegraphics[width=0.84\textwidth]{figures/result_q3_box_delivery.pdf}
\caption{问题三80箱的实际与期望交付时刻；虚线表示两者相等}
\label{fig:q3-box}
\end{figure}

正式验证记录包含36400个运输机采样时刻（分段最大采样间隔1\,s，并补入服务窗边界附近事件点），其中21717个由网关直连，14683个由中继保障，未出现无链路的采样点。连续证书进一步覆盖全部36198.242\,s需求时间：其中直连整段证明21140.352\,s，占58.40\%；中继两跳整段证明15057.890\,s，占41.60\%，未决区间为零。采样点数量与证书区间时长属于不同口径，不能互换。

保障比例在架次间并不平均。\texttt{Q2-009}需要中继证明1206.341\,s，占该架次通信需求时间的59.888\%；\texttt{Q2-012}相应为957.887\,s，占58.845\%。前者又包含全方案最小链路裕量区间，后者承运在硬截止时刻恰好交付的\texttt{S002}两箱。由此可见，单纯统计全局41.60\%的中继服务占比，仍会掩盖关键路线的时间和通信双重紧张程度。此处的占比是证书选择的直连与中继分配比例；若某段同时有两种可用链路，证书只记录其中一种，故不能直接解释为“必须依赖中继”的物理下界。

逐架次返航余量、硬时限、运输机及电池、中继机及能源组件均通过回代；最小运输与中继返航电量比例分别为23.088\%和31.289\%。表~\ref{tab:q3-verification}把“样点无中断”与“全时间可证”分开列出，避免以前者代替后者。运输电耗由问题二的69.580\,kWh变为69.621\,kWh，原因之一是\texttt{Q2-009}调整访问顺序；加上4.215\,kWh中继电耗后为73.835\,kWh。联合完工时刻也较问题二延后1774.815\,s。两问的路线顺序和约束集合不同，这些数值是所得方案的比较，不是通信保障的最小能耗或最短工期证明。
\begin{table}[htbp]
\centering\small
\caption{问题三正式方案的分层可行性核查}
\label{tab:q3-verification}
\begin{tabular}{@{}p{0.20\textwidth}p{0.51\textwidth}p{0.19\textwidth}@{}}
\toprule
核查层 & 核查依据 & 结果\\
\midrule
货箱与硬时限 & 80箱逐箱交付，适用箱用实际时刻比较硬截止 & 漏交/硬违反均为0\\
设备与能源 & 运输机、电池、中继机、组件的具体占用区间及返航电量 & 资源冲突/能量违反均为0\\
事件采样 & 36400个轨迹样点与服务窗事件点的双向链路 & 无链路样点0\\
连续通信 & 1317个区间、36198.242\,s，含双向回传检查 & 未决区间0\\
\bottomrule
\end{tabular}
\end{table}

束宽、候选点位置、固定的初始8个架次和设备编号以及实例时间细化均限制了搜索空间，正式方案是经严格验证的可行启发式解。证书的最小保守链路裕量较小，若实际DEM、定位或传播损耗偏离输入，应重新计算中继位置和服务窗口。

